\documentclass[12pt]{article}
\usepackage[a3paper, margin=1in]{geometry}
\usepackage{amsmath, amssymb, amsthm, ulem, booktabs}

\begin{document}

\section*{Domácí úkol 2.1}

Vypracovali: Daniel \textit{``Localhost"} Ransdorf, Jan \textit{``kokeshh"} Kodeš \hfill Podpis: \rule{4cm}{0.4pt}


\bigskip

Z pozorování z přednášky víme, že pokud je báze B ortonormální, pak pro skalární součin platí:

\[
  \left<\ \vec{u},\vec{v}\ \right> = \left[\ \vec{u}\ \right]_B \cdot \left[\ \vec{v}\ \right]_B
\]

Rozepišme vektory vyjádřene vzhledem k bázi do součinu matice přechodu se samotnými vektory, tím dostaneme:

\[
  \left<\ \vec{u},\vec{v}\ \right> = \left[id\right]_B^K \vec{u} \cdot \left[id\right]_B^K \vec{v}
\]

A tedy nalezněme matici přechodu od kanonické báze k bázi B (inverz k $\left[id\right]_B^K $):

\[
  \left( \left[id\right]_K^B \ |\ I \right) = \left(\begin{array}{cc|cc}
    1&3 & 1&0 \\
    2&7 & 0&1
  \end{array}\right)
  \overset{\text{Gauss--Jordan}}{\sim}
  \left(\begin{array}{cc|cc}
    1&3 & 1&0 \\
    0&1 & -2&1
  \end{array}\right)
  \overset{\text{G--J}}{\sim}
  \left(\begin{array}{cc|cc}
    1&0 & 7&-3 \\
    0&1 & -2&1
  \end{array}\right)
  = \left( I \ |\ \left[id\right]_B^K \right)
\]

Dosaďme do vzorce součinu:

\[
  \left<\left(\begin{array}{c}
    x_1 \\ x_2
  \end{array}\right),\left(\begin{array}{c}
    y_1 \\ y_2
  \end{array}\right)
  \right> = \left(\begin{array}{cc}
    7&-3 \\
    -2&1
  \end{array}\right) \left(\begin{array}{c}
    x_1 \\ x_2
  \end{array}\right)
  \cdot
  \left(\begin{array}{cc}
    7&-3 \\
    -2&1
  \end{array}\right)
  \left(\begin{array}{c}
    y_1 \\ y_2
  \end{array}\right)
  =
  \left(\begin{array}{c}
    7x_1 - 3x_2 \\
    -2x_1 + x_2
  \end{array}\right)
  \cdot
  \left(\begin{array}{c}
    7y_1 - 3y_2 \\
    -2y_1 + y_2
  \end{array}\right)
\]
\[
  (7x_1 - 3x_2)(7y_1 - 3y_2) + (-2x_1 + x_2)(-2y_1 + y_2)
  = 49x_1y_1 -21x_1y_2 - 21x_2y_1 + 9x_2y_2 + 4x_1y_1 - 2x_1y_2 -2x_2y_1 + x_2y_2
\]
\[
  \left<\left(\begin{array}{c}
    x_1 \\ x_2
  \end{array}\right),\left(\begin{array}{c}
    y_1 \\ y_2
  \end{array}\right)
  \right> = \underline{\underline{53x_1y_1 -23x_1y_2 - 23x_2y_1 + 10x_2y_2}}
\]

\end{document}
